\documentclass[Chp3.tex]{subfiles}

\begin{document}

\section{Comparing inequality aversion measures to the racial wealth gap}
\label{subsec:compute-measures}

\subsection{Data}

\par I use the triennial Survey of Consumer Finances (SCF), $1989$--$2022$, and restrict the sample to Black and White households in each wave. For each year, I compute weighted medians for net wealth and gross assets, and then compute $I^{O}$ and $I^{A}$ under two treatments: (i) a restricted sample that keeps only strictly positive wealth values, and (ii) the full sample which adds a year-specific constant so that all observations shifted to be strictly positive.\footnote{Appendix Figure~\ref{fig:shift-mass-diagnostic} plots the shift constant and the weighted share with non-positive wealth.}

\par The \textit{ambiguity} aggregator used for $I^{A}$ follows the Gilboa--Schmeidler construction with $C=\Delta(S)$ and $S=\{\text{Black},\text{White}\}$. Because the criterion is linear in probabilities over the two states, the inner minimum is attained by putting all weight on whichever racial group has the lower within-group expected utility that year, so computationally $W^{A}_{t}=\min\{\mathbb{E}_{B,t}[U(y)],\,\mathbb{E}_{W,t}[U(y)]\}$. The wedge $I^{A}_{t}>I^{O}_{t}$ arises relative to the objective benchmark $W^{O}$ defined in Appendix~\ref{app:objective-benchmark}.

\par With these modeling objects fixed, the theoretical welfare expressions map directly into computable objects in the data: for each year and wealth concept, we compute $W^{O}$ and $W^{A}$, invert $U$ to recover the equally distributed equivalent wealth level, and then construct $I^{O}$ and $I^{A}$. The remaining step is to set values for $\epsilon$, which governs the curvature of $U(y)$ and therefore the degree of inequality aversion imposed in the computation.

\par The parameter $\epsilon$ indexes the curvature of the social welfare function $U(y)$ and therefore the strength of inequality aversion: larger $\epsilon$ means stronger concavity, a higher penalty on low-wealth outcomes, and thus higher implied inequality aversion in both $I^{O}$ and $I^{A}$.

\par For orientation, read the CRRA index in four regions: (i) $\epsilon=0$ is the affine (linear) benchmark with no inequality aversion in the objective aggregator, reported separately in Appendix~\ref{app:eps0-benchmark}; (ii) $0<\epsilon<1$ yields concavity that is still milder than the logarithmic case; (iii) $\epsilon=1$ is the log boundary; (iv) $\epsilon>1$ imposes tail sensitivity strictly stronger than the log case. The subsections below walk through $\epsilon\in\{0.25,0.5,1,1.5,2\}$ so that each reported value sits in a clearly labeled part of this map.

\subsection{The Black-White gap in net wealth and gross assets}

\begin{figure}[H]
\centering
\includegraphics[width=0.92\linewidth]{\subfix{../Figures/03_mean_net_gross_by_race.png}}
\caption{Mean net wealth (left panel) and mean gross assets (right panel) for Black and White households, SCF 1989--2022. Vertical axis: nominal dollars in thousands (so a tick value of $500$ means $\$500{,}000$).}
\label{fig:mean-net-gross-by-race}
\end{figure}

\par Figure~\ref{fig:mean-net-gross-by-race} plots race-specific means before differencing: White household means lie above Black means in both net wealth and gross assets throughout the sample.

\begin{figure}[H]
\centering
\includegraphics[width=0.92\linewidth]{\subfix{../Figures/03_gap_diagnostic_all_obs.png}}
\caption{The Black-White gap in net wealth and gross assets: median wealth differential, SCF 1989--2022.}
\label{fig:gap-net-gross}
\end{figure}

\par Figure~\ref{fig:gap-net-gross} shows that the median Black-White differential is positive and persistent in both series, with gross-asset gaps generally larger than net-wealth gaps over the sample.

\par Checking gross assets alongside net wealth is useful because it verifies that the empirical trend is not an artifact of liability measurement alone: both series move with similar broad dynamics over time. At the same time, the gross-asset gap is higher in levels than the net-wealth gap, which is consistent with definition since net wealth subtracts debts while gross assets do not.

\subsection{Low concavity: $\epsilon=0.25$}

\begin{figure}[H]
\centering
\includegraphics[width=0.92\linewidth]{\subfix{../Figures/08_measures_eps025_net.png}}\\[0.75em]
\includegraphics[width=0.92\linewidth]{\subfix{../Figures/08_measures_eps025_gross.png}}
\caption{Inequality version measures for $\epsilon=0.25$ (top: net wealth, bottom: gross assets, left: shifted sample, right restricted sample).}
\label{fig:eps025-measures}
\end{figure}

\par Even with a low level of inequality aversion (curvature near zero), the ambiguity ranking already yields higher inequality levels than the objective ranking in both wealth concepts (\(I^{A}>I^{O}\)). When I restrict the sample to respondents with positive wealth only, inequality is extremely high and stable across time for the wealth inequality aversion measure. For the full sample (shifted so that everyone has positive wealth), there is more variation over the time period, with some years showing significantly low levels of inequality. The key takeaway is that \textit{the inequality aversion measure reached using the choice under objective uncertainty framework consistently reports that inequality is lower than it is under the ambiguity aversion measure}. 

\subsection{Mild concavity: $\epsilon=0.5$}

\begin{figure}[H]
\centering
\includegraphics[width=0.92\linewidth]{\subfix{../Figures/05_measures_eps05_net.png}}\\[0.75em]
\includegraphics[width=0.92\linewidth]{\subfix{../Figures/05_measures_eps05_gross.png}}
\caption{Inequality aversion measures for $\epsilon=0.5$.}
\label{fig:eps05-measures}
\end{figure}

\par For $\epsilon=0.5$, the performance of the wealth inequality aversion measure $I^{A}$ in the restricted sample is the same: high and comparatively stable inequality levels over time. This time, the inequality aversion measure $I^{O}$ is now higher for this higher value of epsilon. Again, the shifted sample shows more variation across year, while retaining the feature that $I^{A}$ is consistently above $I^{O}$; so the finding that the objective measure understates inequality relative to the ambiguity measure is preserved.

\subsection{Logarithmic case: $\epsilon=1$}

\begin{figure}[H]
\centering
\includegraphics[width=0.92\linewidth]{\subfix{../Figures/06_measures_eps1_net.png}}\\[0.75em]
\includegraphics[width=0.92\linewidth]{\subfix{../Figures/06_measures_eps1_gross.png}}
\caption{Inequality aversion measures for $\epsilon=1$.}
\label{fig:eps1-measures}
\end{figure}

\par The same pattern holds $\epsilon=1$. For both samples,  $I^{A}>I^{O}$, and in the restricted-sample, the objective measure $I^{O}$ has moved even closer to $I^{A}$. An interesting observation throughout these figures is that, for the shifted sample, both measures consistently report higher inequality levels in gross assets than in net wealth, especially between the years 1998-2010. This pattern is robust to values of $\epsilon$.

\subsection{Strong concavity: $\epsilon=1.5$}

\begin{figure}[H]
\centering
\includegraphics[width=0.92\linewidth]{\subfix{../Figures/09_measures_eps15_net.png}}\\[0.75em]
\includegraphics[width=0.92\linewidth]{\subfix{../Figures/09_measures_eps15_gross.png}}
\caption{Inequality aversion measures for $\epsilon=1.5$.}
\label{fig:eps15-measures}
\end{figure}

\par At this level of concavity in $U(y)$, the restricted sample shows that the measures $I^{A}$ and $I^{O}$ track each other very closely. Both are extremely high and stable over time. The pattern persists for the shifted sample, but reported levels of inequality are significantly larger here (for gross assets, there are multiple periods where inequality is essentially $1$.). Still, the finding that the objective uncertainty aversion measure $I^{O}$ underestimates inequality compared to the ambiguity aversion measure $I^{A}$ is preserved. 

\subsection{Strongest curvature: $\epsilon=2$}

\begin{figure}[H]
\centering
\includegraphics[width=0.92\linewidth]{\subfix{../Figures/07_measures_eps2_net.png}}\\[0.75em]
\includegraphics[width=0.92\linewidth]{\subfix{../Figures/07_measures_eps2_gross.png}}
\caption{Inequality aversion measures for $\epsilon=2$.}
\label{fig:eps2-measures}
\end{figure}

\par This is the largest value of $\epsilon$ reported in the main text; it is the most averse to inequality in the sense that the social welfare functional penalizes low wealth sharply.

\par As $\epsilon$ rises from $1.5$ to $2$, the restricted sample essentially shows no difference between $I^{A}$ and $I^{O}$. The variation in inequality over time in the shifted sample is the most pronounced here, and are almost identical for net wealth and gross assets.

\subsubsection{Final takeaways}

\par Taken together, the inequality-measure series tell the same within-sample story across specifications (from $\epsilon=0.25$ through $\epsilon=2$): rising (or at the very least, extremely high) inequality over the time period. Interestingly, when we restrict the sample to positive wealth values, the inequality aversion measures $I^{A}$ and $I^{O}$ are stable over time. However, Figure~\ref{fig:gap-net-gross-restricted} shows that net-wealth and gross-assets median gaps still vary over time and trend upward. This suggests that the theoretical indices are unusually stable when applied to the restricted sample. Additionally, in the restricted sample, the wealth inequality aversion measure $I^{A}$ is especially robust; across reported $\epsilon$ values, it is stable over time and very close to 1. 

\par By contrast, in the shifted sample, treatment both $I^{O}$ and $I^{A}$ are more variable over time and show episodes of comparatively low measured inequality at low and moderate $\epsilon$; they begin to pull upward noticeably at $\epsilon=1.5$, and become frequently near one only in the highest-curvature case $\epsilon=2$. This higher time variation in the shifted series tracks the median-gap dynamics in Figure~\ref{fig:gap-net-gross} much more closely than the restricted-sample series.

\end{document}
